%% cover_letter.tex — A&A submission cover letter
%% Compile: pdflatex cover_letter.tex
\documentclass[12pt,a4paper]{letter}

\usepackage[utf8]{inputenc}
\usepackage[T1]{fontenc}
\usepackage[a4paper, top=2.5cm, bottom=2.5cm, left=2.8cm, right=2.8cm]{geometry}
\usepackage{hyperref}

\signature{Werner Scheibenpflug \\
           Independent Researcher, Vienna, Austria \\
           \texttt{wscheibenpflug@gmail.com} \\
           ORCID: 0009-0008-6450-479X}

\address{Werner Scheibenpflug \\
         Independent Researcher \\
         Vienna, Austria \\
         \texttt{wscheibenpflug@gmail.com}}

\date{\today}

\begin{document}

\begin{letter}{The Editorial Board \\
               Astronomy \& Astrophysics \\
               EDP Sciences}

\opening{Dear Editor,}

I submit for consideration as a Regular Paper in \textit{Astronomy \&
Astrophysics}:

\medskip
\begin{center}
\textbf{``GAPC: A Catalog of HG Phase Curve Parameters for 128,885\\
Main-Belt Asteroids from Gaia DR3 Sparse Photometry''}
\end{center}
\medskip

GAPC is an open-source, reproducible catalog derived from Gaia DR3
sparse photometry, cross-matched with ten external datasets (NEOWISE,
LCDB, Bus-DeMeo/PDS taxonomy, Pravec binaries, SDSS MOC4, DAMIT,
Goffin masses, GASP spectra, AstDys proper elements, MPC). The catalog
contains 128,885 objects and 162 columns; an independent verification
script confirms 75/75 checks pass on the released v8.

\medskip
\textbf{Main scientific results:}

\begin{enumerate}

\item \textbf{Composition-independent size law.} The partial correlation
$r(G,\log D \mid \log p_V) \approx -0.28$ is consistent across all
taxonomy complexes (S, M, E, P, C) and is reproduced within asteroid
families (Flora $n=8{,}168$; Koronis $n=1{,}039$; Eunomia $n=4{,}963$),
pointing to regolith depth as the physical driver rather than
mineralogy.

\item \textbf{Binary $G$-excess.} Known binary asteroids (Pravec
catalog, $n=384$) show a statistically significant $G$-excess after size
control ($p = 2.2 \times 10^{-11}$), providing a new observational
constraint for binary formation and tidal evolution models.

\item \textbf{Phase curves and spectra are orthogonal observables.}
Gaia 16-band reflectance spectra achieve $R^2 \approx 0$ in predicting
$G$ beyond taxonomy and size, confirming that phase curve shape encodes
sub-visual surface texture (grain size, porosity) not accessible through
reflectance spectroscopy.

\end{enumerate}

Several physically motivated correlations --- family age, rotation
period, NEATM thermal inertia, and C-complex hydration --- are found to
be non-significant and are reported explicitly to avoid publication bias.

\medskip
The catalog is available on Zenodo (DOI:
\href{https://doi.org/10.5281/zenodo.19858420}{10.5281/zenodo.19858420})
under MIT license. The full 54-step pipeline is included for
reproducibility.

A companion paper (GASP, Scheibenpflug 2026, \textit{Research Notes of
the AAS}, vol.\ 10, p.\ 87, DOI:
\href{https://doi.org/10.3847/2515-5172/ae5e45}{10.3847/2515-5172/ae5e45})
describing the spectral catalog used in result~(3) has been published.

\medskip
I declare no conflicts of interest. This manuscript has not been
submitted elsewhere and is not under consideration at any other journal.

\closing{Sincerely,}

\end{letter}
\end{document}
